(1) 总自旋 \(\boldsymbol S=\boldsymbol S_1+\boldsymbol S_2+\boldsymbol S_3\) 满足 \(\boldsymbol S^2=\sum_i\boldsymbol S_i^2+2\sum_{i<j}\boldsymbol S_i\cdot\boldsymbol S_j\)。每个自旋半有 \(\boldsymbol S_i^2=3\hbar^2/4\)，故 \(H=\frac J2[S(S+1)-9/4]\)。它与总自旋平方对易，因此按总 \(S\) 标记能级；必须注意三个自旋半张量积包含两个不同的 \(S=1/2\) 不可约表示。

(2) 三个自旋半相加允许一个 \(S=3/2\) 四重态和两个 \(S=1/2\) 双重态，维数检验为 \(4+2+2=8\)。对应能量 \(E_{3/2}=J(15/4-9/4)/2=3J/4\)，\(E_{1/2}=J(3/4-9/4)/2=-3J/4\)。由于两个 \(S=1/2\) 表示能量相同，低能量总简并度为4；未受挫的两自旋反铁磁只有单重态唯一基态。

(3) 正则配分函数 \(Z=4e^{-\beta(3J/4)}+4e^{+\beta(3J/4)}=8\cosh(3\beta J/4)\)。令 \(x=3J/(4k_BT)\)，由 \(U=-\partial_\beta\ln Z\) 得 \(U=-(3J/4)\tanh x\)，热容 \(C=\partial_T U=k_B x^2\operatorname{sech}^2x\)。高温内能趋零，因为完整八维空间上 \(H\) 的迹为零；低温内能趋 \(-3J/4\)。

(4) 熵 \(S=k_B(\ln Z+\beta U)=k_B[\ln(8\cosh x)-x\tanh x]\)。\(T\to0^+\) 时 \(\ln(8\cosh x)\sim x+\ln4\)，所以残余熵 \(S_0=k_B\ln4\)。这是理想孤立等边反铁磁三角形的四重简并与几何阻挫；真实系统可因微弱不等交换、外场或环境耦合解除简并，不能把理论上的残余熵视为必然实验常数。
